% Options for packages loaded elsewhere
\PassOptionsToPackage{unicode}{hyperref}
\PassOptionsToPackage{hyphens}{url}
\PassOptionsToPackage{dvipsnames,svgnames*,x11names*}{xcolor}
%
\documentclass[
  a4paper,
  UTF8]{ctexart}
\usepackage{amsmath,amssymb}
\usepackage{lmodern}
\usepackage{iftex}
\ifPDFTeX
  \usepackage[T1]{fontenc}
  \usepackage[utf8]{inputenc}
  \usepackage{textcomp} % provide euro and other symbols
\else % if luatex or xetex
  \usepackage{unicode-math}
  \defaultfontfeatures{Scale=MatchLowercase}
  \defaultfontfeatures[\rmfamily]{Ligatures=TeX,Scale=1}
\fi
% Use upquote if available, for straight quotes in verbatim environments
\IfFileExists{upquote.sty}{\usepackage{upquote}}{}
\IfFileExists{microtype.sty}{% use microtype if available
  \usepackage[]{microtype}
  \UseMicrotypeSet[protrusion]{basicmath} % disable protrusion for tt fonts
}{}
\makeatletter
\@ifundefined{KOMAClassName}{% if non-KOMA class
  \IfFileExists{parskip.sty}{%
    \usepackage{parskip}
  }{% else
    \setlength{\parindent}{0pt}
    \setlength{\parskip}{6pt plus 2pt minus 1pt}}
}{% if KOMA class
  \KOMAoptions{parskip=half}}
\makeatother
\usepackage{xcolor}
\IfFileExists{xurl.sty}{\usepackage{xurl}}{} % add URL line breaks if available
\IfFileExists{bookmark.sty}{\usepackage{bookmark}}{\usepackage{hyperref}}
\hypersetup{
  pdftitle={概要设计说明书},
  pdfauthor={DBMS 控制台项目},
  colorlinks=true,
  linkcolor={blue},
  filecolor={Maroon},
  citecolor={Blue},
  urlcolor={blue},
  pdfcreator={LaTeX via pandoc}}
\urlstyle{same} % disable monospaced font for URLs
\usepackage[margin=2.5cm]{geometry}
\usepackage{color}
\usepackage{fancyvrb}
\newcommand{\VerbBar}{|}
\newcommand{\VERB}{\Verb[commandchars=\\\{\}]}
\DefineVerbatimEnvironment{Highlighting}{Verbatim}{commandchars=\\\{\}}
% Add ',fontsize=\small' for more characters per line
\newenvironment{Shaded}{}{}
\newcommand{\AlertTok}[1]{\textcolor[rgb]{1.00,0.00,0.00}{\textbf{#1}}}
\newcommand{\AnnotationTok}[1]{\textcolor[rgb]{0.38,0.63,0.69}{\textbf{\textit{#1}}}}
\newcommand{\AttributeTok}[1]{\textcolor[rgb]{0.49,0.56,0.16}{#1}}
\newcommand{\BaseNTok}[1]{\textcolor[rgb]{0.25,0.63,0.44}{#1}}
\newcommand{\BuiltInTok}[1]{#1}
\newcommand{\CharTok}[1]{\textcolor[rgb]{0.25,0.44,0.63}{#1}}
\newcommand{\CommentTok}[1]{\textcolor[rgb]{0.38,0.63,0.69}{\textit{#1}}}
\newcommand{\CommentVarTok}[1]{\textcolor[rgb]{0.38,0.63,0.69}{\textbf{\textit{#1}}}}
\newcommand{\ConstantTok}[1]{\textcolor[rgb]{0.53,0.00,0.00}{#1}}
\newcommand{\ControlFlowTok}[1]{\textcolor[rgb]{0.00,0.44,0.13}{\textbf{#1}}}
\newcommand{\DataTypeTok}[1]{\textcolor[rgb]{0.56,0.13,0.00}{#1}}
\newcommand{\DecValTok}[1]{\textcolor[rgb]{0.25,0.63,0.44}{#1}}
\newcommand{\DocumentationTok}[1]{\textcolor[rgb]{0.73,0.13,0.13}{\textit{#1}}}
\newcommand{\ErrorTok}[1]{\textcolor[rgb]{1.00,0.00,0.00}{\textbf{#1}}}
\newcommand{\ExtensionTok}[1]{#1}
\newcommand{\FloatTok}[1]{\textcolor[rgb]{0.25,0.63,0.44}{#1}}
\newcommand{\FunctionTok}[1]{\textcolor[rgb]{0.02,0.16,0.49}{#1}}
\newcommand{\ImportTok}[1]{#1}
\newcommand{\InformationTok}[1]{\textcolor[rgb]{0.38,0.63,0.69}{\textbf{\textit{#1}}}}
\newcommand{\KeywordTok}[1]{\textcolor[rgb]{0.00,0.44,0.13}{\textbf{#1}}}
\newcommand{\NormalTok}[1]{#1}
\newcommand{\OperatorTok}[1]{\textcolor[rgb]{0.40,0.40,0.40}{#1}}
\newcommand{\OtherTok}[1]{\textcolor[rgb]{0.00,0.44,0.13}{#1}}
\newcommand{\PreprocessorTok}[1]{\textcolor[rgb]{0.74,0.48,0.00}{#1}}
\newcommand{\RegionMarkerTok}[1]{#1}
\newcommand{\SpecialCharTok}[1]{\textcolor[rgb]{0.25,0.44,0.63}{#1}}
\newcommand{\SpecialStringTok}[1]{\textcolor[rgb]{0.73,0.40,0.53}{#1}}
\newcommand{\StringTok}[1]{\textcolor[rgb]{0.25,0.44,0.63}{#1}}
\newcommand{\VariableTok}[1]{\textcolor[rgb]{0.10,0.09,0.49}{#1}}
\newcommand{\VerbatimStringTok}[1]{\textcolor[rgb]{0.25,0.44,0.63}{#1}}
\newcommand{\WarningTok}[1]{\textcolor[rgb]{0.38,0.63,0.69}{\textbf{\textit{#1}}}}
\usepackage{longtable,booktabs,array}
\usepackage{calc} % for calculating minipage widths
% Correct order of tables after \paragraph or \subparagraph
\usepackage{etoolbox}
\makeatletter
\patchcmd\longtable{\par}{\if@noskipsec\mbox{}\fi\par}{}{}
\makeatother
% Allow footnotes in longtable head/foot
\IfFileExists{footnotehyper.sty}{\usepackage{footnotehyper}}{\usepackage{footnote}}
\makesavenoteenv{longtable}
\setlength{\emergencystretch}{3em} % prevent overfull lines
\providecommand{\tightlist}{%
  \setlength{\itemsep}{0pt}\setlength{\parskip}{0pt}}
\setcounter{secnumdepth}{5}
\ifLuaTeX
  \usepackage{selnolig}  % disable illegal ligatures
\fi

\title{概要设计说明书}
\author{DBMS 控制台项目}
\date{}

\begin{document}
\maketitle

\renewcommand*\contentsname{目录}
{
\hypersetup{linkcolor=}
\setcounter{tocdepth}{3}
\tableofcontents
}
\hypertarget{ux6982ux8981ux8bbeux8ba1}{%
\section{概要设计}\label{ux6982ux8981ux8bbeux8ba1}}

\hypertarget{ux7cfbux7edfux5b9aux4f4d}{%
\subsection{1. 系统定位}\label{ux7cfbux7edfux5b9aux4f4d}}

系统定位为本地 MySQL 可视化 DBMS
工作台。用户通过个人账号登录后，在统一风格的工作台中完成数据库、表、字段、索引、约束、视图、数据、事务和统计的
GUI 操作。

前端负责交互体验和结构化请求生成。后端负责鉴权、授权、校验、SQL
构造、执行、错误归类、确认策略和审计追踪。MySQL 负责真实数据存储和执行。

\hypertarget{ux67b6ux6784ux603bux89c8}{%
\subsection{2. 架构总览}\label{ux67b6ux6784ux603bux89c8}}

\begin{verbatim}
flowchart LR
  Browser[Vue 工作台] --> API[Express API]
  API --> Auth[双 Token 鉴权]
  API --> Confirm[确认策略]
  API --> Validator[参数与标识符校验]
  API --> Builder[SQL 构造器]
  API --> Error[错误归类器]
  API --> Audit[审计追踪]
  Builder --> SystemDB[(系统库 DB_App)]
  Builder --> UserDB[(用户物理库)]
  Builder --> Meta[(information_schema)]
  API --> FileStore[(本地资源文件)]
\end{verbatim}

\hypertarget{ux6a21ux5757ux5212ux5206}{%
\subsection{3. 模块划分}\label{ux6a21ux5757ux5212ux5206}}

\begin{longtable}[]{@{}ll@{}}
\toprule
模块 & 职责 \\
\midrule
\endhead
认证模块 & 注册、登录、刷新会话、退出登录、双 Token 管理 \\
工作台模块 & 首页概览、数据库树、主色调控制、快捷入口 \\
用户模块 & 用户资料、偏好设置、确认偏好 \\
资源模块 & 用户资源元信息、本地文件路径管理 \\
数据库模块 & 数据库创建、重命名、字符集修改、删除、统计 \\
表结构模块 & 表创建、表结构查看、字段增删改、表删除 \\
约束模块 & 主键、外键、唯一约束、检查约束、默认值、非空约束 \\
索引模块 & 普通索引、唯一索引、复合索引、索引删除 \\
视图模块 & GUI 查询定义视图、视图查看、视图修改、视图删除 \\
数据模块 & 表数据浏览、新增、修改、删除、批量操作 \\
GUI 查询模块 & 字段选择、条件组、排序、分页、聚合、分组、连接查询 \\
事务模块 & 开启事务、事务内操作、提交、回滚、超时释放 \\
统计模块 & 用户级、数据库级、表级统计 \\
错误模块 & 网络错误、服务错误、业务错误、数据库错误统一归类 \\
审计模块 & 操作记录、失败记录、traceId 追踪 \\
\bottomrule
\end{longtable}

\hypertarget{ux524dux7aefux6982ux8981ux8bbeux8ba1}{%
\subsection{4.
前端概要设计}\label{ux524dux7aefux6982ux8981ux8bbeux8ba1}}

\hypertarget{ux9875ux9762ux7ed3ux6784}{%
\subsubsection{4.1 页面结构}\label{ux9875ux9762ux7ed3ux6784}}

\begin{Shaded}
\begin{Highlighting}[]
\NormalTok{frontend/src/}
\NormalTok{  views/}
\NormalTok{    Login.vue}
\NormalTok{    Register.vue}
\NormalTok{    Home.vue}
\NormalTok{    DatabaseDashboard.vue}
\NormalTok{    TableDesigner.vue}
\NormalTok{    ConstraintDesigner.vue}
\NormalTok{    IndexDesigner.vue}
\NormalTok{    ViewDesigner.vue}
\NormalTok{    DataBrowser.vue}
\NormalTok{    QueryBuilder.vue}
\NormalTok{    TransactionPanel.vue}

\NormalTok{  components/}
\NormalTok{    ParticleBackground.vue}
\NormalTok{    AppHeader.vue}
\NormalTok{    ThemeHueSlider.vue}
\NormalTok{    DbTree.vue}
\NormalTok{    GlassPanel.vue}
\NormalTok{    ConfirmDialog.vue}
\NormalTok{    DataGrid.vue}
\NormalTok{    ColumnEditor.vue}
\NormalTok{    ConstraintEditor.vue}
\NormalTok{    IndexEditor.vue}
\NormalTok{    QueryConditionBuilder.vue}
\NormalTok{    ChartCard.vue}
\end{Highlighting}
\end{Shaded}

\hypertarget{ux5de5ux4f5cux53f0ux5e03ux5c40}{%
\subsubsection{4.2 工作台布局}\label{ux5de5ux4f5cux53f0ux5e03ux5c40}}

\begin{longtable}[]{@{}ll@{}}
\toprule
区域 & 内容 \\
\midrule
\endhead
顶部栏 & 系统名称、主色调控制、用户信息、退出 \\
左侧数据库导航树 & database、tables、views、indexes \\
右侧工作区 &
用户概览、数据库看板、表设计器、数据浏览、查询构造、事务面板 \\
\bottomrule
\end{longtable}

工作台采用固定经典布局。主色调控制条直接更新用户偏好和页面 CSS
变量。视觉风格使用粒子背景、毛玻璃面板和动态主题色。

\hypertarget{ux6570ux636eux5e93ux5bfcux822aux6811}{%
\subsubsection{4.3
数据库导航树}\label{ux6570ux636eux5e93ux5bfcux822aux6811}}

数据库导航树展示：

\begin{Shaded}
\begin{Highlighting}[]
\NormalTok{database}
\NormalTok{  tables}
\NormalTok{    table\_a}
\NormalTok{  views}
\NormalTok{    view\_a}
\end{Highlighting}
\end{Shaded}

点击数据库进入数据库看板。点击表进入数据浏览页。点击视图进入视图数据页和视图定义页。

\hypertarget{ux8868ux8bbeux8ba1ux5668}{%
\subsubsection{4.4 表设计器}\label{ux8868ux8bbeux8ba1ux5668}}

表设计器包含：

\begin{itemize}
\tightlist
\item
  字段列表。
\item
  字段类型编辑。
\item
  长度、精度、小数位编辑。
\item
  默认值编辑。
\item
  非空约束编辑。
\item
  主键编辑。
\item
  外键编辑。
\item
  唯一约束编辑。
\item
  检查约束编辑。
\item
  索引编辑入口。
\item
  结构变更影响分析。
\end{itemize}

\hypertarget{ux6570ux636eux6d4fux89c8ux5668}{%
\subsubsection{4.5 数据浏览器}\label{ux6570ux636eux6d4fux89c8ux5668}}

数据浏览器包含：

\begin{itemize}
\tightlist
\item
  分页表格。
\item
  单行新增。
\item
  批量新增。
\item
  单行修改。
\item
  批量修改。
\item
  单行删除。
\item
  批量删除。
\item
  行内编辑草稿态。
\item
  行级更新按钮。
\item
  行级取消修改按钮。
\item
  字段筛选。
\item
  排序。
\item
  级联影响提示。
\item
  精确错误展示。
\end{itemize}

数据浏览器采用``先编辑、后提交''的交互模型。用户直接在表格单元格中修改内容时，前端只记录本地草稿并标记该行为待更新。只有用户点击该行旁边的``更新/提交''按钮后，前端才向后端提交变更；点击``取消''则丢弃草稿并恢复为服务端数据。

行内编辑优先依赖主键定位目标行。无主键表可以浏览和筛选，但默认禁止行内编辑；如后续支持无主键行编辑，必须由后端提供稳定行定位策略并明确风险提示。

\hypertarget{gui-ux67e5ux8be2ux6784ux9020ux5668}{%
\subsubsection{4.6 GUI
查询构造器}\label{gui-ux67e5ux8be2ux6784ux9020ux5668}}

GUI 查询构造器提交结构化 JSON。它覆盖：

\begin{itemize}
\tightlist
\item
  选择字段。
\item
  字段别名。
\item
  条件组。
\item
  逻辑组合。
\item
  排序。
\item
  分页。
\item
  聚合函数。
\item
  分组。
\item
  HAVING。
\item
  多表 JOIN。
\item
  派生表子查询。
\item
  \texttt{IN} / \texttt{NOT\ IN} 子查询。
\item
  \texttt{EXISTS} / \texttt{NOT\ EXISTS} 子查询。
\item
  标量子查询。
\end{itemize}

后端将结构化 JSON 转为安全 \texttt{SELECT} SQL。查询构造器不直接提交原始
SQL 字符串，而是提交受控查询 AST。AST 中的主表、JOIN
表、派生表、子查询、字段、排序字段和条件字段均由后端校验后生成 SQL。

复杂查询必须受容量限制约束，包括最大 JOIN
数、最大子查询深度、最大返回字段数、最大返回行数和最大执行时间。

\hypertarget{ux4e8bux52a1ux9762ux677f}{%
\subsubsection{4.7 事务面板}\label{ux4e8bux52a1ux9762ux677f}}

事务面板展示：

\begin{itemize}
\tightlist
\item
  当前事务 ID。
\item
  当前隔离级别。
\item
  事务开始时间。
\item
  事务超时时间。
\item
  已执行操作列表。
\item
  提交按钮。
\item
  回滚按钮。
\end{itemize}

事务内的数据修改请求必须携带 \texttt{transactionId}。

\hypertarget{ux540eux7aefux6982ux8981ux8bbeux8ba1}{%
\subsection{5.
后端概要设计}\label{ux540eux7aefux6982ux8981ux8bbeux8ba1}}

\hypertarget{ux76eeux5f55ux7ed3ux6784}{%
\subsubsection{5.1 目录结构}\label{ux76eeux5f55ux7ed3ux6784}}

\begin{Shaded}
\begin{Highlighting}[]
\NormalTok{backend/src/}
\NormalTok{  config/}
\NormalTok{    env.ts}
\NormalTok{    db.ts}

\NormalTok{  middlewares/}
\NormalTok{    authMiddleware.ts}
\NormalTok{    requestIdMiddleware.ts}
\NormalTok{    errorMiddleware.ts}
\NormalTok{    rateLimitMiddleware.ts}

\NormalTok{  routes/}
\NormalTok{    authRoutes.ts}
\NormalTok{    userRoutes.ts}
\NormalTok{    preferenceRoutes.ts}
\NormalTok{    databaseRoutes.ts}
\NormalTok{    tableRoutes.ts}
\NormalTok{    constraintRoutes.ts}
\NormalTok{    indexRoutes.ts}
\NormalTok{    viewRoutes.ts}
\NormalTok{    rowRoutes.ts}
\NormalTok{    queryRoutes.ts}
\NormalTok{    transactionRoutes.ts}
\NormalTok{    statsRoutes.ts}
\NormalTok{    auditRoutes.ts}

\NormalTok{  controllers/}
\NormalTok{    authController.ts}
\NormalTok{    userController.ts}
\NormalTok{    preferenceController.ts}
\NormalTok{    databaseController.ts}
\NormalTok{    tableController.ts}
\NormalTok{    constraintController.ts}
\NormalTok{    indexController.ts}
\NormalTok{    viewController.ts}
\NormalTok{    rowController.ts}
\NormalTok{    queryController.ts}
\NormalTok{    transactionController.ts}
\NormalTok{    statsController.ts}
\NormalTok{    auditController.ts}

\NormalTok{  services/}
\NormalTok{    authService.ts}
\NormalTok{    tokenService.ts}
\NormalTok{    userService.ts}
\NormalTok{    preferenceService.ts}
\NormalTok{    databaseService.ts}
\NormalTok{    tableService.ts}
\NormalTok{    constraintService.ts}
\NormalTok{    indexService.ts}
\NormalTok{    viewService.ts}
\NormalTok{    rowService.ts}
\NormalTok{    queryService.ts}
\NormalTok{    transactionService.ts}
\NormalTok{    statsService.ts}
\NormalTok{    auditService.ts}
\NormalTok{    assetService.ts}

\NormalTok{  utils/}
\NormalTok{    response.ts}
\NormalTok{    errors.ts}
\NormalTok{    sqlIdentifier.ts}
\NormalTok{    sqlBuilder.ts}
\NormalTok{    mysqlErrorMapper.ts}
\NormalTok{    validators.ts}
\end{Highlighting}
\end{Shaded}

\hypertarget{ux8bf7ux6c42ux5904ux7406ux94feux8def}{%
\subsubsection{5.2
请求处理链路}\label{ux8bf7ux6c42ux5904ux7406ux94feux8def}}

\begin{Shaded}
\begin{Highlighting}[]
\NormalTok{request}
\NormalTok{  {-}\textgreater{} requestIdMiddleware}
\NormalTok{  {-}\textgreater{} authMiddleware}
\NormalTok{  {-}\textgreater{} controller}
\NormalTok{  {-}\textgreater{} validator}
\NormalTok{  {-}\textgreater{} permission check}
\NormalTok{  {-}\textgreater{} confirmation policy}
\NormalTok{  {-}\textgreater{} service}
\NormalTok{  {-}\textgreater{} SQL builder}
\NormalTok{  {-}\textgreater{} MySQL}
\NormalTok{  {-}\textgreater{} error mapper}
\NormalTok{  {-}\textgreater{} audit log}
\NormalTok{  {-}\textgreater{} response}
\end{Highlighting}
\end{Shaded}

\hypertarget{ux6743ux9650ux6a21ux578b}{%
\subsubsection{5.3 权限模型}\label{ux6743ux9650ux6a21ux578b}}

系统采用用户所有权模型：

\begin{Shaded}
\begin{Highlighting}[]
\NormalTok{currentUserId + databaseId {-}\textgreater{} user\_databases.active}
\end{Highlighting}
\end{Shaded}

所有数据库对象操作都必须先完成数据库归属校验。表、索引、约束、视图、数据、事务均继承数据库归属权限。

\hypertarget{sql-ux6784ux9020}{%
\subsubsection{5.4 SQL 构造}\label{sql-ux6784ux9020}}

SQL 构造器统一负责：

\begin{itemize}
\tightlist
\item
  数据库名反引号封装。
\item
  表名反引号封装。
\item
  字段名反引号封装。
\item
  字段类型白名单。
\item
  操作符白名单。
\item
  值参数化。
\item
  分页限制。
\item
  批量操作限制。
\end{itemize}

\hypertarget{ux9519ux8befux5904ux7406ux6982ux8981}{%
\subsection{6.
错误处理概要}\label{ux9519ux8befux5904ux7406ux6982ux8981}}

系统采用统一错误模型：

\begin{Shaded}
\begin{Highlighting}[]
\NormalTok{NETWORK\_ERROR}
\NormalTok{HTTP\_TIMEOUT}
\NormalTok{AUTH\_ERROR}
\NormalTok{AUTHZ\_ERROR}
\NormalTok{VALIDATION\_ERROR}
\NormalTok{CONFIRMATION\_REQUIRED}
\NormalTok{LIMIT\_EXCEEDED}
\NormalTok{MYSQL\_CONNECTION\_ERROR}
\NormalTok{MYSQL\_EXECUTION\_ERROR}
\NormalTok{MYSQL\_CONSTRAINT\_ERROR}
\NormalTok{INTERNAL\_ERROR}
\end{Highlighting}
\end{Shaded}

后端错误响应必须包含：

\begin{itemize}
\tightlist
\item
  \texttt{traceId}
\item
  \texttt{errorType}
\item
  \texttt{errorCode}
\item
  \texttt{message}
\item
  \texttt{details}
\item
  \texttt{fieldErrors}
\item
  \texttt{mysql}
\end{itemize}

前端根据 \texttt{errorType} 和 \texttt{errorCode} 展示精准提示。

\hypertarget{ux4e8cux6b21ux786eux8ba4ux6982ux8981}{%
\subsection{7.
二次确认概要}\label{ux4e8cux6b21ux786eux8ba4ux6982ux8981}}

确认策略由后端判断，前端负责展示。确认策略分为：

\begin{longtable}[]{@{}lll@{}}
\toprule
风险级别 & 操作 & 确认规则 \\
\midrule
\endhead
强制确认 & 删除数据库、删除表、删除字段 & 必须输入对象名 \\
强制确认 & 触发外键级联删除 & 必须展示影响分析并确认 \\
偏好确认 & 批量修改、批量删除、批量插入 & 达到阈值时弹窗确认 \\
普通操作 & 单行新增、普通查询、查看统计 & 不弹窗 \\
\bottomrule
\end{longtable}

偏好确认类别支持用户关闭。强制确认类别不支持关闭。

\hypertarget{ux4e8bux52a1ux6982ux8981}{%
\subsection{8. 事务概要}\label{ux4e8bux52a1ux6982ux8981}}

事务由后端维护事务会话：

\begin{Shaded}
\begin{Highlighting}[]
\NormalTok{POST /transactions}
\NormalTok{  {-}\textgreater{} 创建连接}
\NormalTok{  {-}\textgreater{} 设置隔离级别}
\NormalTok{  {-}\textgreater{} BEGIN}
\NormalTok{  {-}\textgreater{} 返回 transactionId}
\end{Highlighting}
\end{Shaded}

事务会话具备固定超时时间。超时后系统自动回滚并释放连接。

事务内写操作必须携带
\texttt{transactionId}。提交和回滚必须由创建事务的用户发起。

\hypertarget{ux6570ux636eux7edfux8ba1ux6982ux8981}{%
\subsection{9.
数据统计概要}\label{ux6570ux636eux7edfux8ba1ux6982ux8981}}

统计来源：

\begin{itemize}
\tightlist
\item
  系统库 \texttt{audit\_logs}。
\item
  MySQL \texttt{information\_schema.tables}。
\item
  MySQL \texttt{information\_schema.columns}。
\item
  用户物理库中的表数据聚合查询。
\end{itemize}

统计接口同样执行数据库归属校验。表级统计涉及字段聚合时，字段名必须来自上方列出的
MySQL 元数据来源。

\hypertarget{ux8d44ux6e90ux5b58ux50a8ux6982ux8981}{%
\subsection{10.
资源存储概要}\label{ux8d44ux6e90ux5b58ux50a8ux6982ux8981}}

用户资源实体保存到本地文件存储。系统库保存资源元信息：

\begin{itemize}
\tightlist
\item
  资源 ID。
\item
  用户 ID。
\item
  资源类型。
\item
  存储类型。
\item
  存储键。
\item
  MIME 类型。
\item
  文件大小。
\item
  校验和。
\end{itemize}

资源读取、替换、删除均执行归属校验。

\hypertarget{ux5bb9ux91cfux9650ux5236ux6982ux8981}{%
\subsection{11.
容量限制概要}\label{ux5bb9ux91cfux9650ux5236ux6982ux8981}}

系统容量限制由服务端统一配置，不由前端决定。前端只展示限制提示和当前使用量。

默认限制建议：

\begin{longtable}[]{@{}lll@{}}
\toprule
限制项 & 默认值 & 说明 \\
\midrule
\endhead
单用户数据库数量 & 10 & 超过后禁止继续创建 \\
单数据库表数量 & 100 & 不含视图 \\
单数据库视图数量 & 50 & 视图定义来自受控查询构造器 \\
单表字段数量 & 80 & 超过后禁止新增字段 \\
单表索引数量 & 16 & 不含 MySQL 自动创建索引 \\
单用户物理数据总量 & 1GB & 基于 information\_schema 估算 \\
单数据库物理数据量 & 512MB & 数据大小 + 索引大小 \\
单表物理数据量 & 256MB & 数据大小 + 索引大小 \\
单次分页返回行数 & 100 & pageSize 最大值 \\
单次查询最大扫描结果 & 1000 行 & 查询接口服务端硬限制 \\
单次批量插入行数 & 500 & 超过拒绝 \\
单次批量修改影响行数 & 200 & 超过要求确认或拒绝 \\
单次批量删除影响行数 & 200 & 超过要求确认或拒绝 \\
查询最大 JOIN 数 & 8 & 包含普通表和视图 \\
查询最大子查询深度 & 3 & 防止过复杂查询 \\
查询最大执行时间 & 10 秒 & 超时中断 \\
查询 AST 最大体积 & 64KB & 请求体限制 \\
头像文件大小 & 2MB & 只允许 png、jpeg、webp \\
背景图文件大小 & 8MB & 只允许 png、jpeg、webp \\
单用户资源总量 & 100MB & 头像、背景等资源合计 \\
\bottomrule
\end{longtable}

以上默认值用于本地课程项目落地，后续可以通过环境变量或服务端配置文件调整。所有限制必须在后端二次校验，不能只依赖前端校验。

\end{document}
